% !TEX root = ../../../main.tex
% Appendix entry: preprocessing.sliding_windows -- "Sliding windows" (Appendix A).
\entryfield{Algorithm}
Everything is in \code{Adapters/preprocessing_sliding_windows.py}; no \code{Working/}
function is called. For a span of $n$ samples at rate $f_s$:
\begin{enumerate}
\item \code{plan_windows}: window length $L = \mathrm{round}(\text{window\_s}\cdot f_s)$
  samples, refused if $L < 2$. The stride is $g = \mathrm{round}(\text{gap\_s}\cdot f_s)$ when
  gap\_s $> 0$, else $g = L$. Despite its name, gap\_s is the distance between window
  \emph{starts}, not the empty time between windows. A stride $g < L$ (overlapping windows)
  is refused. Starts are $0, g, 2g, \dots$ up to the last start that admits a whole window;
  none fitting is an error.
\item \code{assign_split}, rule \texttt{blocked\_by\_time}: block edges
  $e_j = j\,n/B$ for $j = 0..B$ ($B$ = n\_blocks, \code{np.linspace}). A window belongs to the
  block its \emph{start} falls in. With $n_t = \mathrm{round}(B\cdot\text{holdout\_frac})$
  and $n_v = \mathrm{round}(B\cdot\text{validation\_frac})$ (Python's \code{round}, which
  rounds halves to even), the last $n_t$ blocks are test, the $n_v$ before them validation,
  the rest train. Labels: 0 train, 1 validation, 2 test.
\item Rule \texttt{random} (only with allow\_random\_split): one uniform draw $u_i$ per
  window from \code{np.random.default_rng(seed)}; $u_i <$ holdout\_frac is test, the next
  validation\_frac is validation.
\item \code{spacing_check} records the smallest stride and whether it is $\geq L$.
\item The starts are made channel-absolute by adding $\mathrm{round}(t_0 f_s)$, the span's
  offset in the channel.
\end{enumerate}

\entryfield{Assumptions}
\begin{itemize}
\item \emph{Units.} The signal's values are never read, only its length; units do not matter.
\item \emph{Sampling rate.} Lengths are given in seconds and converted with the recording's
  $f_s$, so the same settings mean the same duration at any rate.
\item \emph{Missing values.} NaNs are not checked. A window over a gap is kept, and the
  block after it decides what to do with the missing samples.
\item \emph{Independence.} The blocked split assumes that dependence between windows fades
  within about one window. Nothing enforces an empty buffer at a block border (below).
\end{itemize}

\entryfield{Outputs}
A \texttt{WindowSet} of length $L$, starts channel-absolute, whose feature table holds one
column, \code{split} (0/1/2). \code{split} is on the bookkeeping list
(\code{BOOKKEEPING_COLUMNS} in \code{Working/types/windowset.py}), so
\blockref{catalogue.cluster} and the classifier drop it before treating the table as
features. \code{meta}: \code{length}, \code{gap} (the stride), \code{n_windows},
\code{split_rule}, \code{span_start}, \code{blocks} (index, start, end, role and window
count of every block; empty for a random split), \code{split_counts}, \code{spacing},
\code{seed}. The block also has a pre-run readout (\code{_derive}) showing the window count
and the split counts.

\entryfield{Cost}
$O(n/g)$ time and memory; the signal is not read. Instant at any span length.

\entryfield{Failure modes}
\begin{itemize}
\item \emph{No buffer at block borders.} Windows tile the span, and roles change at a block
  border without any window being dropped, so the last training window can end exactly
  where the first test window begins. Two neighbouring windows still leak through the slow
  autocorrelation at every border. The training pool is stricter: it carves a gap at every
  role change (\blockref{preprocessing.window_pool}), and \code{blocked_split} in
  \code{Working/training/windows.py} drops windows within one window of a border. This
  block does neither.
\item \emph{A window can reach into the next block.} Roles follow the window's start. When
  the block edge $n/B$ is not a multiple of the stride, the last window of a training block
  runs on into the test block's time. It shares no samples with any test window, but it
  does cover test-block time.
\item \emph{Rounding to zero blocks.} With $B = 5$, validation\_frac $= 0.1$ gives
  $\mathrm{round}(0.5) = 0$ validation blocks, and holdout\_frac $= 0.1$ gives no test block.
  Unlike \code{Working/training/windows.py::blocked_split}, this block neither forces at least
  one test block nor refuses a plan that leaves nothing to train on: holdout and validation
  fractions summing to more than one give every block a held-out role.
\item \emph{One window per block or fewer.} Few windows and many blocks leave blocks empty;
  the counts in \code{meta["blocks"]} show it, and the readout warns under ten windows.
\item \emph{Random split.} Allowed only when ticked, and then leaks neighbours by design.
\end{itemize}

\entryfield{Validation}
The leakage argument for blocked splits on dependent data is set out by
\citet{roberts2017cross}. The tests check the geometry and the refusals, not any
downstream effect on a model's score.

\entryfield{References}
\citet{roberts2017cross} on blocked cross-validation. The window sets and their split, as the
Library shows them: \apppage{library.window-sets}.

\entryfield{Tests}
\code{tests/test_adapter_sliding_windows.py}: windows tile the span at gap\_s $= 0$; starts
are channel-absolute; a span shorter than a window and a stride shorter than the window are
refused; the blocked split holds out whole last blocks; a random split is refused unless
accepted; cluster and classifier ignore the \code{split} column.
\code{tests/test_adapter_event_windows.py} checks that \blockref{catalogue.event_windows}
uses the same \code{assign_split}.
